%% The course sequencing chart for the 1000-level algebra and precalculus
%% courses, drawn in TikZ.  Included by \SequencingChart, which osucourse.cls
%% defines; it is a fragment, not a standalone document.
%%
%% To preview it on its own, put it in a standalone document:
%%   \documentclass[tikz,border=8pt]{standalone}
%%   \usetikzlibrary{arrows.meta}
%%   \begin{document}%% The course sequencing chart for the 1000-level algebra and precalculus
%% courses, drawn in TikZ.  Included by \SequencingChart, which osucourse.cls
%% defines; it is a fragment, not a standalone document.
%%
%% To preview it on its own, put it in a standalone document:
%%   \documentclass[tikz,border=8pt]{standalone}
%%   \usetikzlibrary{arrows.meta}
%%   \begin{document}%% The course sequencing chart for the 1000-level algebra and precalculus
%% courses, drawn in TikZ.  Included by \SequencingChart, which osucourse.cls
%% defines; it is a fragment, not a standalone document.
%%
%% To preview it on its own, put it in a standalone document:
%%   \documentclass[tikz,border=8pt]{standalone}
%%   \usetikzlibrary{arrows.meta}
%%   \begin{document}%% The course sequencing chart for the 1000-level algebra and precalculus
%% courses, drawn in TikZ.  Included by \SequencingChart, which osucourse.cls
%% defines; it is a fragment, not a standalone document.
%%
%% To preview it on its own, put it in a standalone document:
%%   \documentclass[tikz,border=8pt]{standalone}
%%   \usetikzlibrary{arrows.meta}
%%   \begin{document}\input{sequencing-chart}\end{document}
\begin{tikzpicture}[
  x=1cm,y=1cm,
  course/.style={draw=black,line width=.65pt,rectangle,
    minimum width=10mm,minimum height=5.8mm,inner sep=2pt,
    font=\sffamily\fontsize{10}{12}\selectfont},
  sequence/.style={-{Triangle[length=1.5mm,width=1.35mm]},line width=.65pt}
]
\node[course] (1050) at (0,0) {1050};
\node[course] (1075) at (1.9,0) {1075};
\node[course] (1116) at (1.9,3) {1116};
\node[course] (1125) at (3.8,3) {1125};
\node[course] (1130) at (3.8,1.5) {1130};
\node[course] (1148) at (3.8,0) {1148};
\node[course] (1131) at (5.7,1.5) {1131};
\node[course] (1149) at (5.7,0) {1149};
\node[course] (1150) at (5.7,-1.5) {1150};
\node[course] (1151) at (8.2,0) {1151};
\node[course] (1156) at (8.2,-1.5) {1156};
\node[course] (1165) at (8.2,-3) {1165};
% Entry and branching from Math 1075.
\draw[sequence] (1050.east) -- (1075.west);
\draw[sequence] (1075.north) -- (1116.south);
\draw[sequence] ([xshift=2mm]1075.north) -- (1125.south west);
\draw[sequence] (1075.north east) -- (1130.south west);
\draw[sequence] (1075.east) -- (1148.west);
% Algebra pathways.
\draw[sequence] (1130.east) -- (1131.west);
\draw[sequence] (1148.north east) -- (1131.south west);
\draw[sequence] (1148.east) -- (1149.west);
% The original chart has no incoming arrow to Math 1150.
\draw[sequence] (1149.east) -- (1151.west);
\draw[sequence] ([yshift=-1mm]1149.east) -- (1156.north west);
\draw[sequence] (1149.south east) -- (1165.north west);
\draw[sequence] (1150.north east) -- (1151.south west);
\draw[sequence] (1150.east) -- (1156.west);
\draw[sequence] (1150.south east) -- (1165.west);
\end{tikzpicture}
\end{document}
\begin{tikzpicture}[
  x=1cm,y=1cm,
  course/.style={draw=black,line width=.65pt,rectangle,
    minimum width=10mm,minimum height=5.8mm,inner sep=2pt,
    font=\sffamily\fontsize{10}{12}\selectfont},
  sequence/.style={-{Triangle[length=1.5mm,width=1.35mm]},line width=.65pt}
]
\node[course] (1050) at (0,0) {1050};
\node[course] (1075) at (1.9,0) {1075};
\node[course] (1116) at (1.9,3) {1116};
\node[course] (1125) at (3.8,3) {1125};
\node[course] (1130) at (3.8,1.5) {1130};
\node[course] (1148) at (3.8,0) {1148};
\node[course] (1131) at (5.7,1.5) {1131};
\node[course] (1149) at (5.7,0) {1149};
\node[course] (1150) at (5.7,-1.5) {1150};
\node[course] (1151) at (8.2,0) {1151};
\node[course] (1156) at (8.2,-1.5) {1156};
\node[course] (1165) at (8.2,-3) {1165};
% Entry and branching from Math 1075.
\draw[sequence] (1050.east) -- (1075.west);
\draw[sequence] (1075.north) -- (1116.south);
\draw[sequence] ([xshift=2mm]1075.north) -- (1125.south west);
\draw[sequence] (1075.north east) -- (1130.south west);
\draw[sequence] (1075.east) -- (1148.west);
% Algebra pathways.
\draw[sequence] (1130.east) -- (1131.west);
\draw[sequence] (1148.north east) -- (1131.south west);
\draw[sequence] (1148.east) -- (1149.west);
% The original chart has no incoming arrow to Math 1150.
\draw[sequence] (1149.east) -- (1151.west);
\draw[sequence] ([yshift=-1mm]1149.east) -- (1156.north west);
\draw[sequence] (1149.south east) -- (1165.north west);
\draw[sequence] (1150.north east) -- (1151.south west);
\draw[sequence] (1150.east) -- (1156.west);
\draw[sequence] (1150.south east) -- (1165.west);
\end{tikzpicture}
\end{document}
\begin{tikzpicture}[
  x=1cm,y=1cm,
  course/.style={draw=black,line width=.65pt,rectangle,
    minimum width=10mm,minimum height=5.8mm,inner sep=2pt,
    font=\sffamily\fontsize{10}{12}\selectfont},
  sequence/.style={-{Triangle[length=1.5mm,width=1.35mm]},line width=.65pt}
]
\node[course] (1050) at (0,0) {1050};
\node[course] (1075) at (1.9,0) {1075};
\node[course] (1116) at (1.9,3) {1116};
\node[course] (1125) at (3.8,3) {1125};
\node[course] (1130) at (3.8,1.5) {1130};
\node[course] (1148) at (3.8,0) {1148};
\node[course] (1131) at (5.7,1.5) {1131};
\node[course] (1149) at (5.7,0) {1149};
\node[course] (1150) at (5.7,-1.5) {1150};
\node[course] (1151) at (8.2,0) {1151};
\node[course] (1156) at (8.2,-1.5) {1156};
\node[course] (1165) at (8.2,-3) {1165};
% Entry and branching from Math 1075.
\draw[sequence] (1050.east) -- (1075.west);
\draw[sequence] (1075.north) -- (1116.south);
\draw[sequence] ([xshift=2mm]1075.north) -- (1125.south west);
\draw[sequence] (1075.north east) -- (1130.south west);
\draw[sequence] (1075.east) -- (1148.west);
% Algebra pathways.
\draw[sequence] (1130.east) -- (1131.west);
\draw[sequence] (1148.north east) -- (1131.south west);
\draw[sequence] (1148.east) -- (1149.west);
% The original chart has no incoming arrow to Math 1150.
\draw[sequence] (1149.east) -- (1151.west);
\draw[sequence] ([yshift=-1mm]1149.east) -- (1156.north west);
\draw[sequence] (1149.south east) -- (1165.north west);
\draw[sequence] (1150.north east) -- (1151.south west);
\draw[sequence] (1150.east) -- (1156.west);
\draw[sequence] (1150.south east) -- (1165.west);
\end{tikzpicture}
\end{document}
\begin{tikzpicture}[
  x=1cm,y=1cm,
  course/.style={draw=black,line width=.65pt,rectangle,
    minimum width=10mm,minimum height=5.8mm,inner sep=2pt,
    font=\sffamily\fontsize{10}{12}\selectfont},
  sequence/.style={-{Triangle[length=1.5mm,width=1.35mm]},line width=.65pt}
]
\node[course] (1050) at (0,0) {1050};
\node[course] (1075) at (1.9,0) {1075};
\node[course] (1116) at (1.9,3) {1116};
\node[course] (1125) at (3.8,3) {1125};
\node[course] (1130) at (3.8,1.5) {1130};
\node[course] (1148) at (3.8,0) {1148};
\node[course] (1131) at (5.7,1.5) {1131};
\node[course] (1149) at (5.7,0) {1149};
\node[course] (1150) at (5.7,-1.5) {1150};
\node[course] (1151) at (8.2,0) {1151};
\node[course] (1156) at (8.2,-1.5) {1156};
\node[course] (1165) at (8.2,-3) {1165};
% Entry and branching from Math 1075.
\draw[sequence] (1050.east) -- (1075.west);
\draw[sequence] (1075.north) -- (1116.south);
\draw[sequence] ([xshift=2mm]1075.north) -- (1125.south west);
\draw[sequence] (1075.north east) -- (1130.south west);
\draw[sequence] (1075.east) -- (1148.west);
% Algebra pathways.
\draw[sequence] (1130.east) -- (1131.west);
\draw[sequence] (1148.north east) -- (1131.south west);
\draw[sequence] (1148.east) -- (1149.west);
% The original chart has no incoming arrow to Math 1150.
\draw[sequence] (1149.east) -- (1151.west);
\draw[sequence] ([yshift=-1mm]1149.east) -- (1156.north west);
\draw[sequence] (1149.south east) -- (1165.north west);
\draw[sequence] (1150.north east) -- (1151.south west);
\draw[sequence] (1150.east) -- (1156.west);
\draw[sequence] (1150.south east) -- (1165.west);
\end{tikzpicture}
